% Compilar con LuaLaTeX. Archivo autónomo; no requiere la clase corporativa.
\documentclass[11pt,letterpaper]{article}
\usepackage[spanish,es-tabla]{babel}
\usepackage{fontspec}
\setmainfont{TeX Gyre Heros}
\usepackage[margin=2cm]{geometry}
\usepackage{longtable,array,booktabs,enumitem,xcolor}
\usepackage[hidelinks]{hyperref}
\usepackage{xurl}
\definecolor{naranja}{HTML}{D75A16}
\setlength{\parindent}{0pt}
\setlength{\parskip}{6pt}
\setlength{\tabcolsep}{4pt}
\renewcommand{\arraystretch}{1.15}
\setcounter{tocdepth}{1}
\setlist[itemize]{leftmargin=*,topsep=2pt,itemsep=2pt,parsep=0pt}
\newcommand{\punteo}[1]{\begin{itemize}[nosep,leftmargin=1em]#1\end{itemize}}
\newcommand{\ruta}[1]{\nolinkurl{#1}}
\newcommand{\inicioTabla}{%
  \small
  \begin{longtable}{@{}>{\raggedright\arraybackslash}p{.13\linewidth}>{\raggedright\arraybackslash}p{.43\linewidth}>{\raggedright\arraybackslash}p{.39\linewidth}@{}}
  \toprule
  \textbf{ID / página de revisión} & \textbf{Punto de la revisión} & \textbf{Cambios realizados que contribuyen a resolverlo} \\
  \midrule\endfirsthead
  \toprule
  \textbf{ID / página de revisión} & \textbf{Punto de la revisión (continuación)} & \textbf{Cambios realizados (continuación)} \\
  \midrule\endhead
  \midrule\multicolumn{3}{r}{\footnotesize Continúa en la página siguiente.}\\\endfoot
  \bottomrule\endlastfoot
}
\newcommand{\finTabla}{\end{longtable}\normalsize}

\title{\textcolor{naranja}{Trazabilidad de la revisión}\\Cambios locales del Subdocumento 7}
\author{Registro de modificaciones y justificación}
\date{5 de noviembre de 2026}

\begin{document}
\maketitle
\tableofcontents
\newpage

\section{Objeto, fuentes y criterio de lectura}

Este registro detalla las modificaciones realizadas en la rama
\ruta{LafroX-Subdoc7-LateX-CORRIGIENDO}. El ámbito editado comprende el
Subdocumento 7 (S7), sus anexos, los formularios T-14, T-15 y T-18, sus figuras,
fuentes y herramientas de regeneración. No se modificaron otros subdocumentos
ni otras ramas.

La fuente evaluadora es \emph{G02 -- LAFROX -- Primera Revisión Informe 2.pdf},
de 55 páginas, ubicado en la carpeta Downloads del usuario. Las páginas de la
columna izquierda remiten a ese PDF. Las páginas que se mencionan como evidencia
de un cambio corresponden a los entregables finales regenerados.

Las tablas incluyen los puntos que requieren corrección o seguimiento en la
revisión, tanto del S7 como del resto del proyecto. Se agrupan las reiteraciones
de una misma observación, incluidas sus repeticiones en «Qué se espera» y en las
consideraciones transversales. Los reconocimientos «OK» y las observaciones del
Informe 1 que el evaluador ya declaró corregidas no se convierten en nuevas
tareas: son antecedentes, no correcciones ejecutadas en esta intervención.

\textbf{Una celda vacía en la columna de cambios significa que en esta
intervención no se realizó ninguna modificación orientada a mejorar ese punto.}
No significa que el punto esté resuelto ni que no aplique. Las celdas con punteo
describen acciones efectivamente realizadas; cuando sólo existe una mejora
parcial, se declara su límite. Los controles y diagnósticos previos, por sí solos,
no se cuentan como correcciones.

Los identificadores de este archivo son propios del registro. No son números
asignados por el evaluador. Se conserva la numeración S7-01 a S7-24 del análisis
previo y se desagregan aspectos adicionales cuando conviene evitar que una mejora
editorial parezca cerrar un problema de planificación.

\section{Cambios realizados: detalle y justificación}

\subsection{C01. Recomposición de las cartas Gantt}
\textbf{Qué cambió.} Se sustituyeron las reproducciones densas por láminas
vectoriales con nombres completos, columnas separadas de inicio y fin, eje
temporal, identificación de filas y leyenda. Los siete tramos del resumen
mantienen sus rangos de filas. Los vínculos que continúan fuera de una lámina se
pueden seguir mediante las hojas de relaciones por ID. También se recompusieron
las vistas de 56 meses, la vista crítica y la presentación de los ocho frentes.

\textbf{Justificación.} Los rótulos pequeños, encabezados superpuestos y códigos
sin leyenda impedían consultar el calendario. Dividir los gráficos mejora su
lectura conservando el contenido de la programación existente.

\textbf{Fuentes y evidencia.} Integración en \ruta{subdocumento-07/contenido.tex},
\ruta{formulario-t14/contenido.tex} y \ruta{formulario-t15/contenido.tex};
herramienta \ruta{revision/gantt_local_vector.py}; activos en
\ruta{imagenes/Diagramas que faltaban/Revision-local/}.
Ejemplo: S7, Figura 7.25, p. 64. Se preservaron los números 7.25 a 7.31.

\textbf{Límite.} Se conservan las 155 filas y 262 relaciones de la fuente del
resumen. No se agregaron predecesoras, no se desagruparon los paquetes del modelo,
no se programaron iteraciones y no se recalculó la ruta crítica. Una mejor
representación de la red existente no demuestra que ésta esté completa.

\subsection{C02. Recomposición de las figuras de EDT}
\textbf{Qué cambió.} Se redibujaron las 17 figuras detalladas y la vista completa
del T-14 en láminas con nombres completos, jerarquía visible y continuaciones.
Se conservó el cuarto nivel de las innovaciones. La comprobación de pertenencia
cubre los 227 paquetes de la vista completa.

\textbf{Justificación.} El texto de las figuras originales resultaba demasiado
pequeño. La tipografía mínima medida de los nuevos gráficos Gantt/EDT es de
9,5 puntos en sus láminas de origen. La división en páginas evita reducir de
nuevo las figuras para hacerlas caber en una única hoja.

\textbf{Fuentes y evidencia.} \ruta{revision/edt_local_vector.py}, integraciones
en S7 y T-14 y activos en la misma carpeta de revisión local. Ejemplo: S7,
Figura 7.10, p. 23, con las agrupaciones y paquetes de innovación.

\textbf{Límite.} Es una mejora propia de presentación. La EDT sigue organizada
por las fases y cuentas existentes: no se ejecutó la reorganización por productos
que solicita el evaluador.

\subsection{C03. Eliminación de una llamada inexistente en la Figura 7.20}
\textbf{Qué cambió.} Se eliminó la indicación «detalle de aceptación H12 abajo»
del objeto de texto del PDF y de la fuente Draw.io recuperada. Se mantuvieron
los hitos y los 33 enlaces. Se ajustó la inclusión de la figura al ancho útil y
se mejoró la identificación de relaciones en el diagrama.

\textbf{Justificación.} La llamada remitía a un detalle que no estaba dibujado.
Su eliminación corrige la inconsistencia sin inventar un procedimiento de
aceptación ni alterar el grafo.

\textbf{Fuentes y evidencia.} Derivados
\ruta{Revision-local/Fig_7-20_Red_Precedencias.drawio} y su PDF;
\ruta{revision/arreglos_locales_sd7.py}. Figura 7.20 del S7, p. 43.
Los activos originales se conservaron.

\subsection{C04. Definición de siglas y leyendas}
\textbf{Qué cambió.} Se incorporaron explicaciones breves en los cinco cuerpos
para SD, EDT, PT, HH, E1/E2, RC, PT/H, H1--H12, F1--F8, QA, SRE, NOC, SOC y DTE.
Las nuevas láminas explican sus códigos y cómo seguir los vínculos.

\textbf{Justificación.} El lector puede interpretar las abreviaturas sin buscar
una definición dispersa en otro formulario. Esta corrección responde a las
observaciones de comprensión, no a un cambio de alcance.

\textbf{Fuentes.} Archivos \ruta{contenido.tex} de S7, anexos, T-14, T-15 y T-18,
y generadores de gráficos locales.

\subsection{C05. Separación entre elaboración y aprobación en H1 y H8}
\textbf{Qué cambió.} En la tabla del diccionario de la cuenta 1.2 se corrigieron
dos períodos:
\begin{itemize}
\item 1.2.1: elaboración preliminar en mes 1; aprobación en H1, mes 2.
\item 1.2.5: elaboración en mes 13; aprobación en H8, mes 14.
\end{itemize}
\textbf{Justificación.} La redacción anterior asociaba el mes de elaboración al
hito de aprobación y parecía contradecir el criterio de aceptación de la misma
fila. La aclaración de H1 responde directamente a la revisión; la de H8 es una
corrección adicional del mismo tipo identificada durante el análisis.

\textbf{Fuente y evidencia.} \ruta{formulario-t14/tablas/tab_T14_4.tex};
T-14, tabla 4, p. 48. No se movieron fechas ni hitos de la programación.

\subsection{C06. Explicación de los objetivos de reversión}
\textbf{Qué cambió.} Se integró junto a la tabla 7.8 del S7 una explicación de
los dos niveles: hasta 10 minutos de reversión técnica; hasta 30 minutos de
preparación operacional en paralelo; y hasta 10 minutos adicionales de
verificación y autorización. El objetivo total se expresa como
\(\max(10,30)+10=40\) minutos. Se mantienen la decisión límite propuesta a las
04:45, la restauración objetivo a las 05:25 y el margen anterior a las 05:30.

Se aclaran el carácter automático del mecanismo técnico y la supervisión SRE,
las responsabilidades del CLIENTE en la contingencia operacional, la medición en
Preproducción y el papel auxiliar del soporte en papel. Se remite a T-18,
sección 6.4, para el procedimiento desarrollado.

\textbf{Justificación.} El cuerpo mencionaba 10 y 40 minutos sin explicar
suficientemente su relación, aunque T-18 ya contenía el desglose.

\textbf{Fuentes y evidencia.} \ruta{subdocumento-07/contenido.tex} y
\ruta{subdocumento-07/tablas/tab_SD7_7.tex}; S7, p. 105.
Son objetivos para un ensayo futuro; no se declara que el ensayo se haya realizado.

\subsection{C07. Simplificación de tablas y conservación de las explicaciones}
\textbf{Qué cambió.} Se acortaron las celdas de cobertura de módulos de la tabla
7.2 y las tres filas de niveles de reversión de la tabla 7.8. Los detalles se
conservaron en prosa cercana, incluida la cobertura del portal del proveedor y
las reglas operacionales de reversión.

\textbf{Justificación.} Las tablas sirven para comparar elementos breves; los
párrafos largos dificultaban seguir columnas y condiciones. No se eliminaron
módulos, cantidades, responsabilidades ni ventanas por abreviar las celdas.

\textbf{Fuentes.} \ruta{subdocumento-07/tablas/tab_SD7_2.tex},
\ruta{subdocumento-07/tablas/tab_SD7_7.tex} y el cuerpo del S7.

\subsection{C08. Depuración parcial de duplicaciones y rótulos}
\textbf{Qué cambió.} Se condensaron los párrafos finales de la sección 7.3.8 del
S7 sobre precio, coordinación de stock y recuperación. Se conservan las reglas
y referencias necesarias para la planificación, remitiendo al detalle ya
existente en T-18, sección 6.5. En éste se sustituyó el título «Controles nuevos
de coordinación y aceptación» por «Verificaciones programadas de coordinación
y recuperación». Se ajustaron expresiones de continuidad entre figuras y sus
leyendas de procedencia.

\textbf{Justificación.} Reducir la repetición facilita distinguir planificación
de diseño técnico. El nuevo título describe el contenido sin presentar el bloque
como una respuesta de una ronda de corrección.

\textbf{Límite.} La observación amplia queda parcialmente atendida: no se migró
todo el contenido técnico a S4 ni se depuró toda la prosa defensiva del proyecto.
Se mantuvieron las advertencias que revelan problemas reales aún pendientes.

\subsection{C09. Comprobación de referencias y composición}
\textbf{Qué cambió.} Se integraron los nuevos activos manteniendo etiquetas y
numeración, se recompilaron los cinco entregables y se verificaron referencias
LaTeX internas, etiquetas duplicadas, glifos ausentes, marcadores de referencia
sin resolver y desbordes. No se detectaron esas incidencias en los archivos finales.

\textbf{Justificación.} Cambiar el número de láminas puede romper llamadas,
índices o numeración. La comprobación evita entregar una fuente actualizada con
un PDF antiguo o con referencias pendientes.

\textbf{Límite.} No certifica la vigencia de todas las fuentes externas ni
equivale a una revisión técnica humana. Persisten avisos heredados de composición
«Infinite glue shrinkage»; no se modificó la clase corporativa para suprimirlos.

\subsection{C10. Regeneración y registro de evidencia}
\textbf{Qué cambió.} Se regeneraron los cinco PDF finales y sus cinco versiones
Markdown. Se sincronizó \ruta{main.pdf} con la compilación del S7. Se adaptó
\ruta{revision/refresh_sd7_markdown.py} para extraer con pypdf y eliminar espacios
finales de línea. Se añadieron controles de rama a las herramientas locales.

Se registraron los hashes iniciales y finales y las comprobaciones en
\ruta{revision/local-20261105/BASE_INICIAL.json},
\ruta{GANTT_VECTOR.json}, \ruta{EDT_VECTOR.json} y \ruta{VERIFICACION.json}.
El informe de cierre está en
\ruta{revision/local-20261105/CIERRE_LOCAL_SD7_2026-11-05.md}.

\textbf{Justificación.} La trazabilidad necesita evidencia reproducible y una
base de comparación. Se comprobó la conservación de las cuatro tablas numéricas
protegidas del T-15 y los tres archivos de programación \ruta{.gan}.

\textbf{Resultado.} Se conservan 227 paquetes, 168 con fecha y 59 continuos,
220.023 HH y un máximo programado de 64 personas. La inspección visual del
asistente cubrió 434 páginas mediante 39 hojas de contacto, con ampliaciones
seleccionadas. No se realizó commit ni push.

\section{Tabla de trazabilidad: Subdocumento 7 y cruces directos}

Las celdas con acciones no implican el cierre completo de la observación. Las
mejoras gráficas conservan el modelo existente; no se atribuyen a las filas de
estimación, dependencias o nivelación porque no se corrigieron esos problemas.

\inicioTabla
S7-01\newline 27--28, 32 & Primer nivel de EDT organizado por fases, sin justificación de agrupación por productos. & \\
S7-02\newline 28, 32 & Etapas como ramas de módulos, pruebas y marcha blanca, en lugar de atributo de los productos. & \\
S7-03\newline 28, 32 & Estimación por clases sin tamaño derivado del T-12 para cada módulo. & \\
S7-04\newline 28, 32 & O/M/P uniformes y probabilidades PERT que no fundamentan la incertidumbre del camino completo. & \\
S7-05\newline 29, 32 & Red incompleta: 102 de 168 paquetes con entregable no declaran predecesora; faltan vínculos de migración, cadenas e innovaciones. & \\
S7-06\newline 29--30, 32 & CPM basado en fechas publicadas, en vez de generar el calendario desde duraciones y relaciones. & \\
S7-07\newline 28--29, 32 & Paralelismo concentrado de módulos y capacidad simultánea de desarrollo sin nivelación demostrada. & \\
S7-08\newline 25, 29--30, 32 & Actividades detalladas sólo hasta H2; no hay iteraciones, ceremonias, demostraciones ni entregas programadas para construcción. & \\
S7-09\newline 25, 30, 40--41 & Pruebas ampliadas concentradas después de construir, sin pruebas visibles por incremento. & \\
S7-10\newline 29, 32 & Dos días de revisión de cierre antes del corte frente al plazo completo de revisión; H7/H12 condicionados. & \\
S7-11\newline 28, 30, 32 & Insuficiencia de calidad: faltan 217,60 HH frente a la dotación de diez personas. & \\
S7-12\newline 28, 30, 32 & Implantación llega a 27 equivalentes sin provisión demostrada; dirección fuera de la tabla corporativa; falta asignación coherente con S1/capítulo 12. & \\
S7-13\newline 36, 38, 53 & Reservas de contingencia, gestión y reserva protegida sin conciliación ni reflejo completo en el calendario. & \\
S7-14\newline 31, 33 & Falta plan para defecto crítico o alto detectado durante el congelamiento de septiembre E2 y su efecto en evidencia y cierre. & \\
S7-15\newline 31 & Inicio tardío del intercambio/certificación de cadenas, sin margen suficiente frente a agosto de 2028. & \\
S7-16\newline 31 & No se concreta la evidencia temprana al proveedor de lácteos entre el inicio y H5. & \\
S7-17\newline 41 & Holgura negativa de H1, H7 y H12 al considerar la subsanación, sin ajuste del cronograma. & \\
S7-18\newline 27 & Contradicción editorial del período de 1.2.1: mes 1 asociado a H1, cuya aprobación es en mes 2. & \punteo{\item C05: separación explícita entre elaboración preliminar del mes 1 y aprobación H1 del mes 2.\item Corrección equivalente adicional de H8: elaboración mes 13 y aprobación mes 14.} \\
S7-19\newline 29--30, 32 & Gantt pequeña; rótulos truncados y encabezados superpuestos. & \punteo{\item C01: láminas vectoriales distribuidas en más páginas; nombres completos y fechas separadas.\item C04: leyendas propias.\item C09: comprobación visual de la composición final.} \\
S7-20\newline 30 & Barras agrupadas dificultan rastrear paquetes y relaciones. & \punteo{\item C01: códigos completos, continuidad por ID y hojas de relaciones para consultar los enlaces de la fuente.\item Avance parcial: se conservan las agrupaciones del modelo; no se demuestra independencia lógica por paquete.} \\
S7-21\newline 32 & Figura 7.20 anuncia un detalle de aceptación H12 que no existe. & \punteo{\item C03: eliminación de la llamada del PDF y de la fuente editable.\item Se conservan los 33 enlaces y los hitos; la red calculada no se reconcilió ni completó.} \\
S7-22\newline 31 & El cuerpo no explica la relación entre los objetivos de reversión de 10 y 40 minutos. & \punteo{\item C06: explicación integrada de niveles, paralelismo y verificación final.\item Fórmula total y horarios explícitos; remisión a T-18, sección 6.4.\item Se distinguen objetivos futuros de resultados de ensayo.} \\
S7-23\newline 31--33 & Tablas 7.2 y 7.8 con párrafos extensos; contenido de diseño en 7.3.8 duplicado en T-18. & \punteo{\item C07: celdas breves y explicaciones en prosa próxima.\item C08: cierre de 7.3.8 condensado y remisión al detalle de T-18.\item Avance parcial: no se trasladó todo el diseño fuera del ámbito S7.} \\
S7-24a\newline 32 & Siglas SD4/SD8, PT/H y RC sin definición suficiente. & \punteo{\item C04: definiciones en los cinco cuerpos y leyendas de gráficos.} \\
S7-24b\newline 32 & Volumen excesivo del cuerpo, formularios y anexos. & \punteo{\item C07--C08: simplificación localizada de celdas y párrafos duplicados.\item Avance limitado: la extensión total aumenta de 316 a 434 páginas por priorizar legibilidad; no se cumple una reducción global.} \\
S7-24c\newline 3, 32, 54 & Foliación reiniciada por archivo; se pide foliación correlativa de la oferta final. & \\
S7-25\newline 32, 54 & Pendiente literal de incompatibilidad reconocido en T-15 y en S7. & \\
S7-26\newline 32 & Prosa defensiva y bloque denominado «Controles nuevos» como respuesta de corrección sin integrar. & \punteo{\item C08: condensación parcial del bloque repetido y nuevo título descriptivo en T-18.\item No se eliminaron advertencias sobre problemas reales ni se depuró toda la prosa defensiva.} \\
S7-27\newline 30 & Remisión a las iteraciones de la sección 6.2, cuya cadencia se especifica en T-10 y no en ese cuerpo. & \\
S7-28\newline 32 & Ruta crítica de construcción E1 que permita sustentar el plazo, además del tramo final. & \\
\finTabla

\section{Tabla de trazabilidad: demás puntos de la revisión}

Este inventario reúne los demás puntos de corrección y seguimiento del documento
evaluador. Sus celdas de cambios se dejan vacías: no se realizaron modificaciones
en esos capítulos. Una mejora del S7 no se registra como solución de una figura,
tabla o procedimiento de otro subdocumento. Las repeticiones transversales se
identifican mediante remisiones para evitar contar dos veces el mismo problema.

\subsection{General y Subdocumento 1}
\inicioTabla
G-01\newline 1--3, 54 & Nomenclatura no conforme del archivo de trazabilidad del Art. 46; integrarlo en una ubicación admitida. & \\
G-02\newline 1--3, 54 & Revisión humana técnica y distribuida, declaración de IA coherente y capacidad de defensa de todos los integrantes. & \\
G-03\newline 2--3 & Reducir volumen general, especialmente S4 y anexos; distinguir síntesis y detalle. La intervención localizada de S7 se registra en S7-24b. & \\
G-04\newline 3 & Erratas de figuras S2: «enERP», «enpapel», «retornarr» y rótulo de diferencias no conciliadas. & \\
G-05\newline 3 & Dominio del sitio corporativo GitHub frente al correo lafrox.cl. & \\
G-06\newline 3 & Páginas oficio de la RBS sin folio/firma y fuentes de tablas S8 con «contenido academico». & \\
S1-01\newline 4--6 & Antecedentes financieros: ventas, patrimonio y capacidad de garantía. & \\
S1-02\newline 5--6 & Una sola versión del SOC: existente, contratado o por implementar. & \\
S1-03\newline 5--6 & Unificar demanda de desarrolladores (31 FTE/40 simultáneos) con T-15 y disponibilidad real. & \\
S1-04\newline 4--6 & Trasladar equipo nominado de 1.5 al capítulo 12; mantener organización del proyecto en S1. & \\
S1-05\newline 5 & Completar organigrama con funciones comerciales/administrativas y ubicación de los dos centros de datos. & \\
S1-06\newline 5--6 & Retirar repetición del plan de renovación de certificaciones. & \\
S1-07\newline 4--6 & Acreditar conocimiento sanitario y de documentos tributarios exigido en RT-15.02. & \\
S1-08\newline 5 & Analizar vencimientos de alianzas entre meses 2 y 11 y su renovación. & \\
S1-09\newline 6 & Explicar referencia farmacéutica con cámaras a menos 22 grados Celsius si se mantiene. & \\
\finTabla

\subsection{Subdocumentos 2 y 3}
\inicioTabla
S2-01\newline 6--8 & Separar problema de solución: promesa 24/48 horas, corte de las 14:00 y retención por SSCC deben tratarse en S3. & \\
S2-02\newline 7--8 & Agregar preparación nocturna, hoja de picking y faltante sin causa al AS-IS de la Figura 2.3. & \\
S2-03\newline 7--8 & Corregir afirmaciones no respaldadas sobre soporte del WMS y módulos satélites del ERP; identificar correctamente la aplicación de preventa. & \\
S2-04\newline 7 & Ubicar rotación del 38 por ciento en el problema de operación de bodega. & \\
S2-05\newline 8 & Completar tabla de restricciones con R-01, R-02, R-03, R-06, R-07, R-09 y R-11. & \\
S2-06\newline 8, 54 & Verificar cita Boulaksil y Belkora: Casablanca en el texto frente a Valencia en la referencia. & \\
S2-07\newline 8 & Corregir declaración de IA de 2.1 que aún atribuye validación de una propuesta de valor. & \\
S3-01\newline 9 & Prioridades con escasa discriminación: 84 por ciento en Crítica o Alta. & \\
S3-02\newline 10--13, 53 & Matriz única de origen, requerimiento, componente, EDT, prueba y criterio, hasta el último eslabón. & \\
S3-03\newline 10 & Desarrollar análisis de interoperabilidad, escalabilidad, desempeño y seguridad más allá de remisiones breves. & \\
S3-04\newline 11, 13 & Rehacer Figura 3.4 legible; representar qué corre en CD, plataformas, camión y nube. & \\
S3-05\newline 11 & Corregir «Gerente de operacions» y rótulos invertidos de las figuras 3.4, 3.5 y 3.7. & \\
S3-06\newline 11 & Recuperar fundamentos de decisiones necesarios para entender Anexo 3.F sin depender de otras lecturas. & \\
S3-07\newline 11, 13 & Completar la columna Origen vacía o cortada en Tabla 3.A.1. & \\
S3-08\newline 11 & Evitar repetir la pila tecnológica completa por fila en T-12; reducir volumen sin perder componente específico. & \\
S3-09\newline 12--13 & Resumir mecanismo de incorporación de conductores externos en 3.4.4. & \\
S3-10\newline 12--13 & Cuantificar dotación de mesa/NOC o remitir expresamente a S10. & \\
S3-11\newline 12--13 & Unificar meta OTIF del cuerpo y Anexo 3.J: mayor que 95 frente a alcanzar 95 por ciento. & \\
S3-12\newline 13, 53 & Usar los mismos nombres de M10 y M11 en S3 y S4. & \\
S3-13\newline 13 & Declaración de IA y revisión técnica coherentes con secciones que contienen figuras. & \\
S3-14\newline 43, 45, 53 & Incorporar las cinco innovaciones al alcance de S3. & \\
\finTabla

\subsection{Subdocumento 4: arquitectura lógica y física}
\inicioTabla
S4-01\newline 14--16 & Figura 1 ilegible; dividir o redibujar con rótulos legibles. & \\
S4-02\newline 14--16 & Trazabilidad requerimiento--módulo--capa por requisito, no sólo por familia; remisión exacta al T-12. & \\
S4-03\newline 15--16 & Leyendas de fuente específicas: indicar los datos que sustentan cada figura. & \\
S4-04\newline 15 & Reducir reproducciones de requisitos que no cambian la arquitectura y el volumen del cuerpo. & \\
S4-05\newline 15--16, 53 & Reconocer y explicar el cambio Django a Laravel respecto del Informe 1. & \\
S4-06\newline 15 & Corregir redacción defectuosa sobre disponibilidad y redundancia. & \\
S4-07\newline 16 & Sustentar existencia y compatibilidad de versiones ofrecidas a la fecha de oferta y tratamiento de dependencias incompatibles. & \\
S4-08\newline 16, 19, 54 & Sustituir referencia inexistente a sección 2.5 del caso por la fuente correcta de clasificación de sala. & \\
S4-09\newline 16, 19 & Revisión humana distribuida y declaración de IA coherente, incluida 4.3 con figuras declaradas como «Ninguno». & \\
S4-10\newline 17, 19, 53 & Plano acotado, superficie y justificación de recintos de la sala técnica. & \\
S4-11\newline 17, 19, 53 & Alternativa si el CLIENTE no aprueba residencia internacional; efectos en latencia, RTO y costo. & \\
S4-12\newline 17, 19, 53 & Estudio de cobertura en cámaras, o método, alcance y fecha de ejecución. & \\
S4-13\newline 18 & Ciclo de vida y plan de reposición por equipo en T-11. & \\
S4-14\newline 18--19 & Figura 18 de topología AWS legible. & \\
S4-15\newline 19 & Cuantificar y llevar al plan de riesgos la residencia exterior y su riesgo regulatorio/latencia. & \\
S4-16\newline 19, 25, 53 & Alinear provisión de cadena de frío por LafroX con exclusión del capítulo 11 del caso. & \\
S4-17\newline 19 & Explicar en Tabla 14 el total por máximo coincidente. & \\
\finTabla

\subsection{Subdocumento 5}
\inicioTabla
S5-01\newline 21--22 & Modelo general Figura 5.1 y vistas por dominio legibles al tamaño de impresión; evitar título duplicado dentro de imagen y leyenda. & \\
S5-02\newline 21 & Unificar significado de M: exclusivamente módulo frente a módulo/mes en el glosario. & \\
S5-03\newline 22 & Fundamentar umbral de desviación monetaria de inventario de 0,3 por ciento. & \\
S5-04\newline 22 & Mostrar cálculo de volúmenes de registros por bytes sin obligar a reconstruirlo en S4. & \\
S5-05\newline 22, 53 & Explicar prosa comprimida, siglas, remisiones y relevos preinscritos de 8/16 horas. & \\
S5-06\newline 22 & Revisión humana técnica del modelo de custodia y de Tabla 5.1. & \\
S5-07\newline 22, 53 & Explicar credencial de turno de 14 horas ante rutas de hasta 14,5 horas y su relación con identidad sin enlace. & \\
\finTabla

\subsection{Subdocumento 6}
\inicioTabla
S6-01\newline 23, 26, 54 & Completar regla truncada de valor ganado: umbral, acción, autoridad y plazo. & \\
S6-02\newline 23--26 & Explicar adaptación PMBOK/RUP, alternativas, convivencia Kanban--iteraciones y límites de adaptación respecto de hitos fijos. & \\
S6-03\newline 23, 26 & Incorporar reglas Art. 72: análisis de impacto, tope acumulado y cambios en etapas de fecha fija. & \\
S6-04\newline 23, 26 & Escalamiento con disparador, destinatario y plazo. & \\
S6-05\newline 23, 26 & Comunicaciones específicas con terreno, turno nocturno, terceros y planificador. & \\
S6-06\newline 23--26, 53 & Unificar 17/19 interesados y completar influencia, interés y estrategia, incluidas empresas transportistas. & \\
S6-07\newline 24, 26, 54 & Retirar duplicación de evaluadores externos entre S6 y T-9 y ubicar su dimensionamiento donde corresponde. & \\
S6-08\newline 24 & Método de seguimiento y nivelación de capacidad propia en solapamientos. & \\
S6-09\newline 24 & Desarrollar gestión de requerimientos, arquitectura evolutiva, refactorización y deuda técnica. & \\
S6-10\newline 24, 26 & Adaptar desarrollo y validación al caso: desconexión, frío, ERP sin documentación, cadenas, planificador y transferencia a cuatro personas de TI. & \\
S6-11\newline 25--26 & Nombrar ambientes, artefactos promovidos y aprobaciones en T-10. & \\
S6-12\newline 25--26 & Programar iteraciones, ceremonias, entregas y pruebas por incremento en Gantt. Corresponde también a S7-08 y S7-09. & \\
S6-13\newline 25--26 & Corregir fuente de Figura 6.1; agregar diagramas de gobierno, iteración y cadena de integración. & \\
S6-14\newline 26 & T-9 y T-10 deben detallar operación, matrices, plantillas, umbrales y plazos, sin repetir el cuerpo. & \\
S6-15\newline 26 & Declaración y revisión humana de contenido coherentes con la autoría efectiva. & \\
\finTabla

\subsection{Subdocumento 8}
\inicioTabla
S8-01\newline 33, 38 & Escala de impacto anclada en plazo, HH, alcance, calidad sanitaria y contrato; definir factores decimales en el cuerpo. & \\
S8-02\newline 33, 38 & Eliminar solape de bandas de probabilidad en 70 por ciento. & \\
S8-03\newline 33 & Conciliar frase sobre zona de prioridad con umbrales de exposición de la tabla. & \\
S8-04\newline 33--34 & Explicar y conciliar probabilidades cualitativas y cuantitativas distintas para el mismo evento. & \\
S8-05\newline 34, 37, 54 & Resolver en diseño/calidad los 13 riesgos internos del proveedor; no trasladarlos al CLIENTE, especialmente los altos/críticos. & \\
S8-06\newline 34, 37, 54 & Identificar contraparte del CLIENTE, solicitud concreta, fecha y contingencia por falta de respuesta. & \\
S8-07\newline 34 & Riesgo de jubilación del planificador con fecha y pérdida de cumplimiento, no sólo indisponibilidad genérica. & \\
S8-08\newline 34, 37 & Incorporar objeción sindical a cámaras y GPS. & \\
S8-09\newline 34 & Incorporar rechazo de conductores de terceros al uso de la solución. & \\
S8-10\newline 35, 37 & Incorporar rotación del 38 por ciento en preparación nocturna. & \\
S8-11\newline 35, 37 & Incorporar interfaces ERP/WMS sin documentación. & \\
S8-12\newline 35, 37 & Tratar enlace único de Concepción de forma específica. & \\
S8-13\newline 35, 37 & Incorporar peak de septiembre sobre la fase del proyecto. & \\
S8-14\newline 35, 38, 53 & Aplicar modos de falla, árbol de fallas o simulación a los riesgos pertinentes. & \\
S8-15\newline 35, 37 & Incorporar obsolescencia, bloqueo por proveedor y escalabilidad. & \\
S8-16\newline 35--37 & Llevar al registro puntos únicos de falla de S4, manifiesto de identidad, cobertura en cámaras y enlaces de cross-docking. & \\
S8-17\newline 35 & Revisar pertinencia de eventos operacionales reescritos sólo como afectación de una prueba frente a riesgos de proyecto ausentes. & \\
S8-18\newline 35--37, 54 & Estimar riesgo residual y plan adicional para todo residual alto/crítico; contingencias específicas. & \\
S8-19\newline 36, 38, 53 & Análisis costo-beneficio de respuestas, corregir antecedente de retorno 0,84 y justificar retornos menores que uno. & \\
S8-20\newline 36--38 & Acciones con esfuerzo, plazo, fecha/hito y contingencia ejecutable. & \\
S8-21\newline 36, 38, 53 & Derivar 1.600 HH de gestión y conciliar reservas de S8 y S7 en T-15 y cronograma. Véase S7-13. & \\
S8-22\newline 36 & Definir y calcular VME mencionado en T-16. & \\
S8-23\newline 36 & Conciliar ventana meses 11--15 de R-034 citada por S13 con Anexo 8.A sin meses. & \\
S8-24\newline 37--38 & Renumerar tablas 8.4 duplicadas y revisar exceso de columnas de Tabla 8.5. & \\
S8-25\newline 37--38, 54 & Quitar «contenido academico», corregir fuentes normativas, tildes y erratas. & \\
S8-26\newline 37 & Incorporar folio/firma a RBS oficio y revisar suficiencia del desarrollo del capítulo. & \\
S8-27\newline 37, 54 & Integrar pendientes declarados, depurar estilo defensivo y efectuar revisión humana técnica de riesgos y T-16. & \\
\finTabla

\subsection{Subdocumentos 9 y 13}
\inicioTabla
S9-01\newline 38--39, 42 & Evidencia de acreditación CMMI vigente o presentación como plan, sin invocar garantía no acreditada. & \\
S9-02\newline 40--42, 53 & Prueba del manifiesto de identidad en corte: emisión, relevo, expiración, revocación y criterio de salida. & \\
S9-03\newline 40, 42 & Usabilidad con recién ingresados al turno nocturno, población explícita y productividad de primera semana. & \\
S9-04\newline 40, 42 & Ambiente de restauraciones mensuales durante septiembre y diciembre. & \\
S9-05\newline 40--41, 54 & Pruebas por iteración y cronograma ajustado por subsanación. Véanse S7-09 y S7-17. & \\
S9-06\newline 41 & Resumen por módulo al inicio de la matriz de casos T-13. & \\
S9-07\newline 41--42 & Corregir errata «extbfCorrecciones en congelamiento» y superposición de «Recuperación» en Figura 9.4. & \\
S9-08\newline 41--42 & Unificar «paso a producción»: apertura a marcha blanca frente a producción desde meses 16 y 21. & \\
S9-09\newline 41--42 & Una sola regla de correcciones durante congelamiento, coherente con S6 y Bases. & \\
S9-10\newline 41--42 & Depurar salvedades defensivas que oscurecen las reglas. & \\
S13-01\newline 43, 45, 53 & Innovaciones explícitas en alcance S3. Véase S3-14. & \\
S13-02\newline 43, 45 & Defender INN-02 como innovación de proceso del CLIENTE o reformularla. & \\
S13-03\newline 43--45, 53 & Resolver doble uso de INN-03 como innovación y cumplimiento de RT-05.30. & \\
S13-04\newline 44--45, 53 & Cuantificar beneficio de cada innovación en unidades físicas con datos del caso. & \\
S13-05\newline 44 & Verificar fechas y ventanas de paquetes con S7, T-14, T-15, T-16 y Anexo 8.A. & \\
S13-06\newline 44, 54 & Sustituir referencia a «Informe 3» por destino propio de la oferta. & \\
S13-07\newline 44--45 & Declaración de IA y revisión técnica efectiva de modelos, umbrales y cálculos. & \\
\finTabla

\section{Conservación de datos y aspectos que no se deben dar por resueltos}

La conservación de cantidades y fechas fue una condición de esta intervención
editorial. No constituye una validación de suficiencia de la planificación.
Permanecen los problemas de estimación, red incompleta, CPM, nivelación, revisión
de actas, congelamientos y reservas identificados en las celdas vacías.

No se modificaron la clase corporativa, los metadatos comunes, la fecha de
portada ni la declaración humana de IA. La foliación integrada pertenece a la
consolidación de la oferta completa. Tampoco se incorporaron contrataciones,
acuerdos del CLIENTE, ensayos ejecutados o resultados de pruebas que no existen.

\begin{center}
\begin{tabular}{lrr}
\toprule
\textbf{Entregable} & \textbf{Páginas antes} & \textbf{Páginas después}\\
\midrule
S7 & 66 & 117\\
Anexos S7 & 25 & 26\\
T-14 & 102 & 166\\
T-15 & 99 & 100\\
T-18 & 24 & 25\\
\midrule
Total & 316 & 434\\
\bottomrule
\end{tabular}
\end{center}

La mayor extensión responde a la división de gráficos para permitir su lectura.
No se presenta como una reducción global de volumen. El presente registro es un
documento de trabajo de trazabilidad; no se incorpora automáticamente a la oferta
ni pretende resolver su nomenclatura contractual.

\end{document}
